%=====================================================================
\section*{Problem 3. Mixture of experts}

Every MoE run replaces each of the eight FFNs of d8w512 by four experts of the same shape with a linear router over the residual stream (\texttt{p3\_moe.py}). Top-1 routing executes one expert per token, so the FFN FLOPs equal the dense model's while the FFN parameters are $4\times$ (105M total). ``Unit'' routing is top-1 selection with the combine weight fixed to 1 instead of the softmax probability: the same expert choice, no router gradient. The diagnostics (load CV across the four experts, routing entropy normalized by $\log 4$, mean largest probability) are measured after training on 8 held-out sequences of 128 tokens; we report the worst layer for the CV and the mean over layers otherwise. Tables~\ref{tab:p3central} to \ref{tab:p3own} give every variant; Figure~\ref{fig:p3curves} the curves.

\begin{table}[H]\centering\small
\caption{Problem 3(a): routing variants (d8w512, LR 0.003). Load CV is the worst layer; entropy and mean max $p$ are means over the MoE layers, on 8 held-out sequences.}\label{tab:p3central}
\begin{tabular}{lrrrr}\toprule
Variant & val loss & max layer load CV & norm. entropy & mean max $p$ \\ \midrule
dense FFN & 2.9267 & -- & -- & -- \\
random top-1 & 3.0840 & 0.084 & -- & -- \\
frozen top-1 & 2.8918 & 1.732 & 0.870 & 0.471 \\
learned top-1 & 2.9067 & 1.732 & 0.617 & 0.654 \\
unit top-1 (matched combine) & 2.9519 & 0.823 & 0.997 & 0.279 \\
all four experts & 2.8271 & 0.000 & 0.783 & 0.537 \\
\bottomrule\end{tabular}
\end{table}
\begin{table}[H]\centering\small
\caption{Problem 3(b): balancing and router conditioning added to learned top-1 routing. $^\dagger$: checkpoint NLL (W\&B history lost).}\label{tab:p3tricks}
\begin{tabular}{lrrrr}\toprule
Variant & val loss & max layer load CV & norm. entropy & mean max $p$ \\ \midrule
balance .01 & 2.8927 & 0.134 & 0.784 & 0.521 \\
balance .001$^\dagger$ & 2.8705 & 0.572 & 0.784 & 0.525 \\
balance .1$^\dagger$ & 2.8964 & 0.165 & 0.818 & 0.483 \\
normalized router & 2.9082 & 1.732 & 0.530 & 0.706 \\
noise .05 & 2.8864 & 1.732 & 0.591 & 0.673 \\
noise .2$^\dagger$ & 2.9040 & 1.732 & 0.422 & 0.778 \\
shared init & 2.9054 & 1.732 & 0.633 & 0.647 \\
quantile routing & 2.8923 & 0.175 & 0.725 & 0.590 \\
norm + balance .01$^\dagger$ & 2.9004 & 0.168 & 0.738 & 0.548 \\
balance .01 + noise .05 & 2.8872 & 0.152 & 0.758 & 0.548 \\
shared init + balance .01 & 2.8851 & 0.129 & 0.783 & 0.526 \\
norm+balance+noise+shared init & 2.8832 & 0.175 & 0.587 & 0.664 \\
\bottomrule\end{tabular}
\end{table}
\begin{table}[H]\centering\small
\caption{Problem 3(c): our router variants; the control is learned top-1 with balance 0 (2.9067).}\label{tab:p3own}
\begin{tabular}{lrrrr}\toprule
Variant & val loss & max layer load CV & norm. entropy & mean max $p$ \\ \midrule
hash router (own) & 2.9312 & 0.141 & 0.949 & 0.365 \\
straight-through router (own)$^\dagger$ & 2.9384 & 0.886 & 0.528 & 0.691 \\
router z-loss 1e-3 (own) & 2.9003 & 1.732 & 0.777 & 0.535 \\
\bottomrule\end{tabular}
\end{table}

\begin{figure}[H]\centering
\includegraphics[width=\linewidth]{p3_curves}
\caption{Validation loss of the routing variants (left) and of the balancing and conditioning variants of learned top-1 routing (right).}
\label{fig:p3curves}
\end{figure}

\subsection*{Part (a). Simple routing}

\begin{itemize}
\item \lead{Random top-1 is far worse than dense} (3.084 against 2.927). The expert is re-drawn per token every step, so each expert is a quarter-size FFN trained on a quarter of the data with no consistency.
\item \lead{A frozen random router beats dense} by 0.035 (2.892). The assignment is a fixed function of the residual stream, so similar hidden states always go to the same expert and the $4\times$ parameters are usable even though nothing was learned about the routing.
\item \lead{Learned top-1 beats dense by 0.020 but loses to frozen by 0.015} (2.907). Its curve shows a loss spike of 0.1 at update 7,200 from which it never fully recovers, the signature of late routing flips. LR 0.003 is the right rate for it: 0.0015 gives 2.914 and 0.006 gives 2.952. A seed-1 rerun gives 2.900 with the same first-layer collapse, so the 0.015 gap to the frozen router is within seed noise; what is robust is that learning the router does not beat a fixed random one.
\item \lead{The matched control says the softmax weight matters twice.} Unit routing keeps the learned argmax but drops the probability from the combine, so the router gets no gradient. It reaches 2.952, worse than dense by 0.025 and than learned top-1 by 0.045. The weight is the router's only learning signal and also scales each expert's output. Unit routing is the most balanced top-1 variant (CV 0.82, entropy 1.00), which shows that balance alone is not the objective.
\item \lead{All four experts set the ceiling.} Executing every expert with its softmax weight ($4\times$ the FFN FLOPs, 1.9$\times$ the model FLOPs) gives 2.827, 0.10 better than dense and 0.08 better than the best top-1 model. Its mean largest probability is only 0.54, so the full model uses two or three experts per token; the top-1 models leave that capacity unused. On the Problem 1 yardstick 2.827 is the ladder's loss at 300 PFLOP, so $\epsilon\approx1.0$: the full mixture is as efficient as a bigger dense model, no better.
\end{itemize}

\lead{What the diagnostics add.} Both the frozen and the learned router are degenerate in the first MoE layer (load CV $\sqrt3=1.73$: every held-out token goes to one expert, so that layer is a dense FFN). In the other seven layers the learned router is more confident (mean max $p$ 0.65 against 0.47, entropy 0.62 against 0.87) and less balanced (per-layer CV 0.55 to 0.67 against near-uniform loads). Training wants a consistent, roughly uniform partition; the learned router buys confidence at the price of balance and stability.

\subsection*{Part (b). Improving the router}

\begin{itemize}
\item \lead{Load balancing fixes the collapse.} Coefficient 0.01 brings the worst-layer CV from 1.73 to 0.13 and the loss from 2.907 to 2.893, level with the frozen router; the spike at 7,200 disappears. Coefficient 0.1 gives 2.896 (CV 0.17). Coefficient 0.001 gives 2.871$^\dagger$, the best top-1 loss we measured, while only half-fixing the imbalance (CV 0.57). So a weak penalty is enough and 0.01 may already over-constrain the partition.
\item \lead{Normalized routing, noise and shared init do not fix the collapse on their own.} With a normalized router (2.908), $\sigma=0.05$ noise (2.886) or cloned experts (2.905) the first MoE layer still sends every token to one expert, and the remaining layers get more confident (mean max $p$ 0.71, 0.67, 0.65). Noise is the only one that helps, by 0.02, by acting as a soft balancing pressure in the surviving layers; normalization costs 0.002.
\item \lead{More noise is not better.} $\sigma=0.1$ gives 2.886, the same as 0.05, and $\sigma=0.2$ gives 2.904$^\dagger$ with mean max $p$ 0.78.
\item \lead{Compositions mostly do not add up.} Normalized router plus balance 0.01 gives 2.900$^\dagger$ (CV 0.17), the balance result and nothing from the norm. All four tricks together (norm, balance 0.01, noise 0.05, shared init) give 2.883 with CV 0.18, about 0.01 better than balance alone, which is the size of the noise effect. Balance 0.01 plus noise 0.05 alone gives 2.887, the same 0.006 gain over balance alone, so noise accounts for all of the ``all tricks'' improvement. Shared init plus balance 0.01 gives 2.885 (CV 0.13), 0.008 better than balance alone, although shared init alone did nothing: once the partition is balanced, starting the experts from one function seems to help them specialize. %COMPOS%
\item \lead{Quantile routing} dispatches by per-expert score quantile with a bias updated without gradients, so every expert gets a quarter of the tokens by construction. It reaches 2.892 with CV 0.17, the same as balance 0.01 and the frozen router. %QNOISE%
\end{itemize}

Balance matters for the loss (random 3.084 and dense 2.927 against balanced 2.893), but through which partition it produces, not through the statistic itself: the unit-combine model is balanced and bad, the frozen model is collapsed in one layer and good. Any reasonable balanced, consistent partition lands at about 2.89 for top-1 at this scale.

\subsection*{Part (c). Our own routers}

Each variant keeps the learned-router model unchanged (four experts, top-1, the probability still scales the expert) and changes one thing; the matched control is learned top-1 with balance 0.
\begin{itemize}
\item \lead{Hash router.} The expert is $\mathrm{id}\bmod 4$ of the input token (Roller et al.), so the partition is content-free, parameter-free and balanced by construction. Prediction: between random (3.084) and learned (2.907), closer to learned, because a consistent assignment lets each expert specialize on a stable slice of the vocabulary. Result: 2.931, CV 0.14, entropy 0.95. Consistency alone is worth 0.15 over random re-draws, and learning which tokens share an expert is worth a further 0.024 (learned) to 0.04 (balanced learned).
\item \lead{Straight-through router.} Hard one-hot combine in the forward pass with the softmax gradient in the backward pass, so each expert's output is not scaled down by $p_e<1$. Prediction: within 0.01 of learned top-1, because the $4p_e$ combine already rescales the output. Result: 2.938$^\dagger$ on the diagnostic set, where learned top-1 scores 2.908, so it is 0.03 worse, outside our prediction. No layer collapses (worst CV 0.89 against 1.73), but the router is the most confident of all variants (mean max $p$ 0.69) with per-layer CVs of 0.16 to 0.89. Passing the full expert output regardless of confidence removes the soft down-weighting of uncertain tokens that the probability-scaled combine provides, which seems to matter more than the extra gradient.
\item \lead{Router z-loss.} Add $10^{-3}\,\mathbb{E}[(\log\sum_e e^{s_e})^2]$ to keep router logits small (ST-MoE). Prediction: fewer late routing flips and a loss close to balance 0.01 if the spike at 7,200 is a logit-scale problem, otherwise no change. Result: 2.900, 0.007 better than learned top-1 (within MoE seed noise), with the first and last MoE layers still collapsed (CV 1.73) and a less confident router (mean max $p$ 0.54 against 0.65). The z-loss tames the logit scale but not the collapse, so it is not a substitute for balancing.
\end{itemize}

\lead{Synthesis.} At this scale almost all of the MoE gain comes from a consistent partition of the tokens: a frozen random projection beats dense by 0.035 and a learned router adds nothing over it unless its collapse is fixed. A small balancing loss (0.001 to 0.01) is the one intervention that matters; normalization, noise and shared initialization change the loss by at most 0.02 and do not compose. Executing all four experts is worth another 0.07 but costs $1.9\times$ the FLOPs, which the Problem 1 yardstick says is a fair trade and not a bargain.

\lead{Rules of thumb.}
\begin{itemize}
\item Four-expert top-1 MoE at d8: $\Delta L$ from dense is $+0.16$ random, $-0.035$ frozen, $-0.020$ learned, $+0.025$ unit combine, $-0.10$ all four experts.
\item The learned router collapses one layer onto one expert (CV 1.73) without balancing; balance 0.01 gives CV 0.13 and $\Delta L=-0.014$; balance 0.001 gives the best top-1 loss, 2.871.
\item Noise 0.05 is worth $-0.02$ on its own; normalization, shared init and router tricks stacked on balance are worth at most $-0.01$ more.
\item A hash router is 0.024 worse than a learned one and 0.15 better than random routing; a straight-through combine is 0.03 worse; a z-loss changes nothing beyond noise.
\end{itemize}
